% TEMPLATE-MILC-v3.tex - plantilla maestra UNICA de las guias MILC v3.
%
% La usan las tres familias de guias, cada una con su driver:
%   · Grados 6-11  -> scripts/build-guias-g11.py
%   · Bebras       -> scripts/build-guias-bebras.py
%   · Semillero    -> scripts/build-guias-semillero.py
%
% Antes habia tres copias casi identicas (~400 lineas iguales, 6 distintas):
% cada mejora habia que aplicarla tres veces, y en la practica no se hacia
% (el motor de recursos vivia solo en dos de ellas). Lo que varia entre
% familias esta parametrizado en 6 placeholders que cada driver llena
% (se nombran aqui SIN su sintaxis real: el builder reemplaza por texto plano,
% tambien dentro de los comentarios, y un valor multilinea rompe el preambulo):
%
%   HEADER_IZQ        encabezado izquierdo de pagina
%   HEADER_DER        encabezado derecho de pagina
%   PORTADA_ETIQUETA  rotulo sobre el numero grande (GRADO/MOMENTO/MODULO)
%   PORTADA_SERIE     linea de serie de la portada
%   PORTADA_ID        identificador de la guia en la portada
%   PORTADA_CIRCULO   el circulito de grado (°), o vacio si no aplica
%
% El resto de claves las documenta content/guias/_SCHEMA.md. El script valida
% que no queden placeholders sin reemplazar antes de escribir el archivo final.

\documentclass[11pt,letterpaper]{article}
\usepackage[margin=2.0cm]{geometry}
\usepackage[table]{xcolor}
\usepackage{fontspec}
\usepackage[spanish]{babel}
\usepackage{tikz}
% Librerías TikZ para los diagramas de las guías (bloque `recursos:`):
%  · babel  → babel en español vuelve activos caracteres como «>», y TikZ los usa
%             en sus opciones (>=stealth, ->). Sin esta librería, cualquier
%             diagrama con flechas revienta con «\language@active@arg> has an
%             extra }». Debe ir ANTES de que se dibuje nada.
%  · positioning        → sintaxis «below left=2pt and 3pt».
%  · decorations.pathreplacing → llaves ({brace}) para señalar rangos.
\usetikzlibrary{babel,positioning,decorations.pathreplacing}
\usepackage{graphicx}
% Circuitos: el trabajo de electrónica se hace hoy con micro:bit y sus sensores
% integrados (MakeCode, sin cableado), así que no se necesitan esquemáticos.
% Si algún día entran montajes con protoboard, basta descomentar esta línea:
% los diagramas .tex/.tikz declarados en `recursos:` ya se incrustan solos.
% \usepackage{circuitikz}
\usepackage[most]{tcolorbox}
\usepackage{tabularx}
\usepackage{array}
\usepackage{fancyhdr}
\usepackage{titlesec}
\usepackage[document]{ragged2e}  % bandera a la derecha en todo el cuerpo
\usepackage{enumitem}
\usepackage{microtype}

% Helvetica Neue no trae la flecha U+2192, asi que cada «→» escrito literalmente
% en el YAML se compilaba a NADA, en silencio (462 flechas en 61 guias: rutas del
% tipo «paleta → tipografia → lockup» salian rotas). Mapeamos el caracter a su
% version matematica, que si tiene glifo. Asi el autor puede seguir escribiendo
% «→» comodamente en el YAML.
\usepackage{newunicodechar}
\newunicodechar{→}{\ensuremath{\rightarrow}}
% Mismo problema con los simbolos de marca y vinetas: el «✓» de «una palomita ✓»
% salia como cuadrito vacio en el PDF de todas las guias que lo usaban.
\usepackage{pifont}
\newunicodechar{✓}{\ding{51}}
\newunicodechar{✗}{\ding{55}}
\newunicodechar{✔}{\ding{52}}
\newunicodechar{•}{\textbullet}
\newunicodechar{≥}{\ensuremath{\geq}}
\newunicodechar{≤}{\ensuremath{\leq}}

\setmainfont{Helvetica Neue}
\setsansfont{Helvetica Neue}
\setlength{\parindent}{0pt}
\setlength{\parskip}{6pt}
% Sin particion de palabras: en 6.º-7.º una palabra cortada por guion al final
% de linea («Comuni-cacion») frena la lectura mucho mas que una linea corta.
\hyphenpenalty=10000
\exhyphenpenalty=10000
\pretolerance=2000
\tolerance=3000
\setlength{\emergencystretch}{3em}
\linespread{1.10}
\renewcommand{\arraystretch}{1.25}
\setlist{leftmargin=5mm,itemsep=1mm,topsep=1mm}
\newcolumntype{Y}{>{\RaggedRight\arraybackslash}X}

\definecolor{milcMagenta}{HTML}{E6007E}
\definecolor{milcVino}{HTML}{5A0038}
\definecolor{milcTurquesa}{HTML}{0093A5}
\definecolor{milcVerde}{HTML}{58A923}
\definecolor{milcMostaza}{HTML}{E5B400}
\definecolor{milcOcre}{HTML}{B86600}
\definecolor{milcNegro}{HTML}{241816}
\definecolor{milcGris}{HTML}{F7F7F4}
\definecolor{milcLinea}{HTML}{E8E2DD}
\definecolor{milcGradePrimary}{HTML}{7C3AED}
\definecolor{milcGradeDark}{HTML}{4C1D95}
\definecolor{milcGradeSoft}{HTML}{EDE9FE}
\definecolor{milcGradeAccent}{HTML}{CFE7FF}
\definecolor{milcGradeText}{HTML}{FFFFFF}
\color{milcNegro}

\pagestyle{fancy}
\fancyhf{}
\lhead{\small Tecnología e Informática · Grado 9}
\rhead{\small Guía 19 · MILC v3}
\cfoot{\small\thepage}
\renewcommand{\headrulewidth}{0pt}

\newtcolorbox{titlebox}[1]{
  enhanced,colback=#1,colframe=#1,arc=7mm,boxrule=0pt,
  left=7mm,right=7mm,top=3mm,bottom=3mm,
  before skip=8pt,after skip=8pt,
  fontupper=\sffamily\bfseries\Large\color{white}
}

\newtcolorbox{softbox}[3]{
  enhanced,colback=#2,colframe=#1,borderline west={4pt}{0pt}{#1},
  arc=2mm,boxrule=.4pt,left=4mm,right=4mm,top=3mm,bottom=3mm,
  before skip=6pt,after skip=8pt,title={#3},coltitle=white,
  colbacktitle=#1,fonttitle=\sffamily\bfseries,fontupper=\RaggedRight,
  attach boxed title to top left={xshift=3mm,yshift=-2mm},
  boxed title style={arc=2mm, boxrule=0pt}
}

\newtcolorbox{drawbox}[2]{
  enhanced,colback=white,colframe=#1,arc=4mm,boxrule=1pt,
  left=5mm,right=5mm,top=4mm,bottom=4mm,before skip=8pt,after skip=8pt,
  title={#2},coltitle=white,colbacktitle=#1,fonttitle=\sffamily\bfseries,
  fontupper=\RaggedRight,
  attach boxed title to top left={xshift=4mm,yshift=-2mm},
  boxed title style={arc=3mm, boxrule=0pt}
}

\newtcolorbox{infoband}[1]{
  enhanced,colback=#1!8,colframe=#1!50,arc=2mm,boxrule=.5pt,
  left=4mm,right=4mm,top=2mm,bottom=2mm,before skip=4pt,after skip=4pt,
  fontupper=\small\RaggedRight
}

% Puente narrativo entre secciones (contrato editorial Conéctate).
% Da continuidad: "Acabas de X. Ahora vas a Y porque Z."
% Visualmente: banda izquierda en color del grado, texto en itálica pequeña.
\newtcolorbox{puentebox}{
  enhanced,colback=milcGradeSoft,colframe=milcGradeDark,
  borderline west={3pt}{0pt}{milcGradeDark},
  arc=0mm,boxrule=0pt,left=5mm,right=4mm,top=2.5mm,bottom=2.5mm,
  before skip=4pt,after skip=8pt,
  fontupper=\itshape\small\color{milcNegro}\RaggedRight
}
% La flecha va en modo matematico a proposito: Helvetica Neue NO tiene el glifo
% U+2192, asi que \textrightarrow se compilaba a NADA («Missing character: There
% is no → in font Helvetica Neue») y el puente salia sin flecha. En modo
% matematico la flecha viene de la fuente matematica y si se dibuja.
\newcommand{\puente}[1]{%
  \begin{puentebox}%
    \textcolor{milcGradeDark}{$\rightarrow$}\hspace{2mm}#1%
  \end{puentebox}%
}

\newcommand{\linead}[1]{\noindent\color{milcLinea}\rule{#1}{0.5pt}\color{milcNegro}}
\newcommand{\respuesta}[1]{\par\vspace{2mm}\foreach \n in {1,...,#1}{\noindent\color{milcLinea}\rule{\linewidth}{0.5pt}\par\vspace{5mm}}\color{milcNegro}}
\newcommand{\checkbox}{\textcolor{milcGradeDark}{\fbox{\phantom{x}}}\hspace{5pt}}

\newcommand{\phasecard}[4]{
\begin{tcolorbox}[enhanced,colback=#3,colframe=#2,arc=3mm,boxrule=.5pt,height=3.05cm,valign=center,halign=center,left=1.5mm,right=1.5mm,top=2mm,bottom=2mm]
{\sffamily\bfseries\fontsize{22}{24}\selectfont\textcolor{#2}{#1}}\par
{\sffamily\bfseries\small #4}
\end{tcolorbox}
}

% ── Piezas del rediseno 2026 (legibilidad 6.º-11.º) ───────────────────
% Banda de estacion: reemplaza el titlebox macizo. Titulo a la izquierda,
% dato practico (tiempo/modalidad) a la derecha, en una sola linea.
\newcommand{\milcestacion}[2]{%
\begin{tcolorbox}[enhanced,colback=#1,colframe=#1,arc=2mm,boxrule=0pt,
  left=5mm,right=5mm,top=2.6mm,bottom=2.6mm,before skip=11pt,after skip=7pt]
{\sffamily\bfseries\fontsize{15}{18}\selectfont\color{white}#2}
\end{tcolorbox}}

% Tarjeta de paso numerada: sustituye las tablas «Pilar / Aplicacion», que
% eran formato de consulta docente, no de estudiante. El numero va en un
% circulo sobre el borde para que la secuencia se lea de un vistazo.
\newcommand{\milcpaso}[3]{%
\begin{tcolorbox}[enhanced,colback=white,colframe=milcLinea,arc=1.5mm,boxrule=.9pt,
  left=14mm,right=4mm,top=3mm,bottom=3mm,before skip=4pt,after skip=4pt,
  fontupper=\RaggedRight,
  overlay={\node[anchor=north west,circle,fill=#2,text=white,inner sep=0pt,
    minimum size=7.4mm,font=\sffamily\bfseries\small]
    at ([xshift=3.6mm,yshift=-3.2mm]frame.north west) {#1};}]
#3
\end{tcolorbox}}

% Casilla marcable de tamano util con lapiz (la anterior era \fbox{\phantom{x}},
% de alto variable segun el texto que la seguia).
\newcommand{\milcmarca}{\framebox[3.8mm]{\rule{0pt}{3.4mm}}}

% Fila marcable de lista suelta (checks de la estacion 1).
\newcommand{\milccheckrow}[1]{%
\par\vspace{1.8mm}\noindent
\makebox[6.4mm][l]{\raisebox{-.55mm}{\textcolor{milcNegro}{\milcmarca}}}%
\parbox[t]{\dimexpr\linewidth-6.4mm\relax}{\RaggedRight #1}\par}

% Celda marcable dentro de la rubrica (tabla de autoevaluacion). Misma
% sangria francesa, pero pensada para vivir dentro de una columna X.
\newcommand{\milcopcion}[1]{%
\makebox[5.2mm][l]{\raisebox{-.55mm}{\textcolor{milcNegro}{\milcmarca}}}%
\parbox[t]{\dimexpr\linewidth-5.2mm\relax}{\RaggedRight #1}}

% ── Recursos visuales de la guía ──────────────────────────────────────
% Declarados en el bloque `recursos:` del YAML y emitidos por el builder.
% \guiaFigura   → imagen rasterizada o PDF (foto, carta celeste, montaje).
% \guiaDiagrama → fuente TikZ/circuitikz (.tex) incrustada como vector:
%                 es la vía para circuitos y esquemáticos editables.
% #1 = ruta relativa al .tex · #2 = pie (puede ir vacío)
\newcommand{\guiaFigura}[2]{%
\begin{center}
\includegraphics[width=0.88\linewidth,height=0.38\textheight,keepaspectratio]{#1}\\[1.5mm]
{\footnotesize\itshape\textcolor{milcNegro!75}{#2}}
\end{center}
}

\newcommand{\guiaDiagrama}[2]{%
\begin{center}
\input{#1}\\[1.5mm]
{\footnotesize\itshape\textcolor{milcNegro!75}{#2}}
\end{center}
}

\begin{document}
\thispagestyle{empty}

% ── Portada ───────────────────────────────────────────────────────────
% Antes: un rectangulo de color a pagina completa. Se veia bien en pantalla,
% pero en la fotocopiadora del colegio se comia el toner y salia gris sucio.
% Ahora la banda ocupa el tercio superior y el resto va en blanco.
\begin{tikzpicture}[remember picture,overlay]
  \path[fill=milcGradePrimary]
    (current page.north west) rectangle ([yshift=-10.6cm]current page.north east);

  \node[anchor=north west,text=milcGradeText] at ([xshift=2.0cm,yshift=-1.55cm]current page.north west)
    {\sffamily\bfseries\fontsize{12}{14}\selectfont GRADO};
  \node[anchor=north west,text=milcGradeText,opacity=.85] at ([xshift=2.0cm,yshift=-2.35cm]current page.north west)
    {\sffamily\bfseries\fontsize{8.5}{10}\selectfont SERIE GUÍAS · TECNOLOGÍA E INFORMÁTICA};
  \node[anchor=north east,text=milcGradeText,opacity=.85,align=right] at ([xshift=-2.0cm,yshift=-1.55cm]current page.north east)
    {\sffamily\bfseries\fontsize{9}{12}\selectfont I.E. Sor María Juliana\\Cartago, Valle del Cauca};

  \node[anchor=west,text=milcGradeText] (gradonum) at ([xshift=1.85cm,yshift=-7.1cm]current page.north west)
    {\sffamily\bfseries\fontsize{112}{112}\selectfont 9};
  % El circulito de grado (°) solo aplica a las guias de grado («11°»); en
  % Bebras y Semillero el numero grande es un momento/modulo, no un grado.
  % Se posiciona relativo al nodo `gradonum`, no en coordenadas absolutas.
  \draw[milcGradeText,line width=1.6pt]
    ([xshift=2.0mm,yshift=-4.5mm]gradonum.north east) circle (.15cm);

  \node[anchor=west,text=milcGradeText,text width=11.4cm,align=left] at ([xshift=1.5cm,yshift=0.5cm]gradonum.east)
    {
     {\sffamily\bfseries\fontsize{9}{11}\selectfont\MakeUppercase{Período 2 · Diseño editorial digital}}\\[3mm]
     {\sffamily\bfseries\fontsize{21}{25}\selectfont\setlength{\baselineskip}{27pt}Maquetar la\\[4pt]revista completa ---\\[4pt]flujo de trabajo}\\[3.5mm]
     {\sffamily\bfseries\fontsize{15}{18}\selectfont Guía 19-9-TIC}
    };
\end{tikzpicture}

\vspace*{9.4cm}
\vfill

{\sffamily\bfseries\fontsize{8}{10}\selectfont\textcolor{milcGradeDark}{LO QUE TE LLEVAS HOY}}

\begin{tcolorbox}[enhanced,colback=white,colframe=milcNegro,arc=2mm,boxrule=1.6pt,
  left=5mm,right=5mm,top=4mm,bottom=4mm,before skip=3pt,after skip=12pt,
  fontupper=\RaggedRight\fontsize{11}{15.5}\selectfont]
Revista digital de 8 páginas con tema propio del estudiante (cultura local, ciencia, deporte, arte, etc.), exportada como PDF y subida a una URL pública (Google Drive, Issuu, Notion o equivalente)
\end{tcolorbox}

\begin{tcolorbox}[enhanced,colback=milcGris,colframe=milcGris,arc=2mm,boxrule=0pt,
  left=5mm,right=5mm,top=3.5mm,bottom=3.5mm,after skip=12pt]
{\sffamily\bfseries\fontsize{8}{10}\selectfont\textcolor{milcNegro!55}{ESTA GUÍA ES DE}}\\[3mm]
\begin{tabularx}{\linewidth}{X p{3.2cm} p{3.2cm}}
\linead{\linewidth} & \linead{3.0cm} & \linead{3.0cm}\\[-1mm]
{\fontsize{8.5}{10}\selectfont\textcolor{milcNegro!55}{Nombre completo}} &
{\fontsize{8.5}{10}\selectfont\textcolor{milcNegro!55}{Grupo}} &
{\fontsize{8.5}{10}\selectfont\textcolor{milcNegro!55}{Fecha}}\\
\end{tabularx}
\end{tcolorbox}

\vfill

\noindent\textcolor{milcLinea}{\rule{\linewidth}{1pt}}\\[2mm]
{\sffamily\bfseries\fontsize{9.5}{12}\selectfont Autor: Dr. Álvaro Cárdenas Orozco}\\[1mm]
{\sffamily\fontsize{8.5}{11}\selectfont\textcolor{milcNegro!55}{Metodología MILC · apertura ancestral · pensamiento computacional · triángulo Dussel–estoicismo–Floridi}}

\newpage

\section*{Apertura: Maquetar la revista completa --- flujo de trabajo}

\begin{softbox}{milcGradeDark}{milcGradeSoft}{Saber ancestral}
Las \textbf{tejedoras Misak del Cauca} arman sus \emph{chumbes} con una planificación previa rigurosa: deciden el patrón general antes de tejer la primera fila, porque cada figura geométrica debe encajar con la siguiente y con el todo. Si una sola fila se equivoca, el calendario codificado se rompe. Esa misma lógica de \emph{ver el conjunto antes de tejer la parte} es la que practicaba el \textbf{diagramador tradicional} en las imprentas del Valle (Cali, Cartago, Tuluá): armaba la revista físicamente, pegando recortes en una mesa, decidiendo qué iba antes y qué después con \emph{todo a la vista}. Los maestros diagramadores tenían mesas enormes donde extendían las 16, 24 o 32 páginas de la revista a la vez. Caminaban alrededor de la mesa, miraban desde lejos, ajustaban. La revista nacía como cuerpo completo, no como páginas sueltas. \textbf{Maquetar digitalmente conserva ese gesto}: tener todas las páginas visibles a la vez para decidir el \emph{ritmo del conjunto}, no solo la belleza de cada parte.
\end{softbox}

\begin{softbox}{milcTurquesa}{milcGris}{Saber contemporáneo}
Hoy maquetar una revista digital de 8 páginas en Figma, Canva o Adobe Express significa: (a) tener \textbf{vista general} de las 8 páginas; (b) verificar que la \textbf{secuencia narrativa} tenga ritmo (no 8 páginas con la misma estructura); (c) mantener \textbf{coherencia del sistema visual} (grilla, tipografía, paleta) a través de todas; (d) \textbf{exportar} como PDF profesional; (e) \textbf{publicar} con URL pública para compartir. Es la sesión de \emph{cosecha} del periodo: aquí converge todo lo aprendido en las 8 sesiones anteriores.
\end{softbox}

\begin{softbox}{milcMostaza}{milcGris}{Pregunta-puente}
\textit{Cuando miras una revista bien hecha y la pasas página por página, ¿qué hace que el conjunto se sienta ``un cuerpo'' y no ``8 páginas pegadas''? Hoy vas a aprender a producir esa coherencia.}
\end{softbox}

\begin{softbox}{milcVerde}{milcGris}{Saber y saber hacer}
Al terminar podrás: (1) \textbf{identificar} los elementos de una revista de 8 páginas y su orden típico; (2) \textbf{analizar} el flujo narrativo y visual de una revista entre páginas; (3) \textbf{aplicar} sistema visual coherente a 8 páginas distintas; (4) \textbf{crear} y publicar tu propia revista digital de 8 páginas con tema propio.
\end{softbox}



\small
\begin{tabularx}{\textwidth}{>{\bfseries}p{3.0cm}Y}
\rowcolor{milcGradePrimary!12}Autor & Dr. Álvaro Cárdenas Orozco\\
DBA & Producción de piezas editoriales digitales con criterio estético y ético (MEN, T\&I)\\
Referentes & Tipografía colombiana · Diseño centrado en accesibilidad · Floridi\\
Producto final & Revista digital de 8 páginas con tema propio del estudiante (cultura local, ciencia, deporte, arte, etc.), exportada como PDF y subida a una URL pública (Google Drive, Issuu, Notion o equivalente)\\
\end{tabularx}
\normalsize

\puente{Antes de empezar, mira el viaje completo. Esta es la sesión más larga e intensa del periodo: en una sesión vas a maquetar 8 páginas. Las sesiones 1-8 te prepararon. Hoy converge todo en producto real.}

\milcestacion{milcGradeDark}{Ruta MILC: así vamos a aprender}
\begin{tabularx}{\textwidth}{>{\centering\arraybackslash}X>{\centering\arraybackslash}X>{\centering\arraybackslash}X>{\centering\arraybackslash}X}
\phasecard{1}{milcVerde}{milcGris}{Escucha\\[-1mm]\small Observo, escucho y pregunto.} &
\phasecard{2}{milcTurquesa}{milcGris}{Entender\\[-1mm]\small Sistematizo y ordeno.} &
\phasecard{3}{milcMagenta}{milcGris}{Hacer\\[-1mm]\small Construyo y aplico.} &
\phasecard{4}{milcVino}{milcGris}{Cierre\\[-1mm]\small Evalúo y reflexiono.}
\end{tabularx}


\puente{Antes de maquetar tu revista, vamos a estudiar la estructura típica de una revista de 8 páginas: qué va en portada, qué en contraportada, cómo se distribuye el contenido en las 6 páginas intermedias. Eso evita la tentación del primer borrador.}

\milcestacion{milcVerde}{Estación 1 de 3 · Escucha}

\begin{softbox}{milcVerde}{milcGris}{Escena para escuchar}
\textbf{Actividad 1 · IDENTIFICA --- Estructura de una revista de 8 páginas} (15 min · individual). Busca \textbf{2 revistas} reales (impresas o digitales) de 8-16 páginas: una cultural seria (Bocas, Cromos, El Malpensante) y una temática (música, deporte, ciencia, naturaleza). Para cada una: ¿qué va en portada? ¿qué en contraportada? ¿cómo se distribuye el contenido principal? ¿hay sumario o índice? ¿hay sección editorial?
\end{softbox}

{\sffamily\bfseries\fontsize{8}{10}\selectfont\textcolor{milcNegro!55}{MARCA CUANDO LO HAYAS HECHO}}
\milccheckrow{Analicé 2 revistas reales con 8+ páginas.}
\milccheckrow{Identifiqué qué va en portada, contraportada y páginas internas.}
\milccheckrow{Reconocí la presencia de sumario, editorial u otras secciones recurrentes.}

\begin{infoband}{milcVerde}


\textbf{Lo que va al cuaderno (Actividad 1 · IDENTIFICA).} Título: «Actividad 1 --- Anatomía de una revista». \textbf{Formato:} esquema dibujado de las 8 páginas de UNA de las 2 revistas analizadas (portada, p.2, p.3...p.8) con etiqueta de qué contiene cada una. \textbf{Sabes que terminaste cuando} puedes nombrar la función de cada una de las 8 páginas en un esquema.
\end{infoband}

\puente{Ya viste la estructura. Ahora vamos a entender el \textbf{flujo de trabajo del diagramador}: 5 pasos para maquetar 8 páginas sin perderte. Sin método, te vas a dispersar a la hora 1 de un trabajo de 4 horas.}

\milcestacion{milcTurquesa}{Estación 2 de 3 · Entender}

\begin{softbox}{milcTurquesa}{milcGris}{Lo que vas a entender}
Maquetar 8 páginas requiere flujo de trabajo, no improvisación. Vamos a descomponerlo con los pilares del pensamiento computacional para que tengas mapa antes de empezar 4 horas de trabajo.
\end{softbox}

\milcpaso{1}{milcTurquesa}{Toda revista de 8 páginas se descompone en estructura típica: (1) \textbf{Portada} (titular principal + imagen impactante + número/fecha), (2) \textbf{Página 2}: editorial o sumario, (3) \textbf{Páginas 3-7}: contenido principal (puede ser 2-3 artículos), (4) \textbf{Página 8}: contraportada (cierre visual, anuncio, crédito o cita destacada). Variaciones existen, pero esa base sirve.}
\milcpaso{2}{milcTurquesa}{Todo flujo de trabajo de maquetación sigue el patrón \textbf{``planear, maquetar, revisar, publicar''}: 30 min planear contenido y secuencia, 2-3 horas maquetar, 30 min revisar/corregir, 30 min exportar y publicar. Saltar planear es la causa principal de revistas que ``se sienten desordenadas''.}
\milcpaso{3}{milcTurquesa}{Lo esencial cabe en una pregunta: \emph{¿al pasar las 8 páginas en orden, hay ritmo o todas se sienten iguales?}. Una revista debe tener variación visual (algunas páginas con más texto, otras con más imagen, otras de respiro) sin perder coherencia.}
\milcpaso{4}{milcTurquesa}{Pasos para maquetar 8 páginas: (1) lista las 8 páginas y qué contendrá cada una (planeación); (2) abre 8 lienzos en tu herramienta; (3) aplica grilla y paleta a todas (coherencia); (4) maqueta una a una, empezando por la más simple; (5) revisa el conjunto cada 2-3 páginas; (6) corrige; (7) exporta como PDF; (8) sube a Drive/Issuu/Notion para URL pública.}

\begin{drawbox}{milcTurquesa}{Anatomía de una revista de 8 páginas}
\textbf{Portada (p.1):} titular + imagen + número/fecha. \\[1mm] \textbf{Editorial/Sumario (p.2):} introducción + lista de contenidos. \\[1mm] \textbf{Artículo principal (p.3-5):} 2-3 páginas, spread completo. \\[1mm] \textbf{Artículo secundario (p.6-7):} 1-2 páginas. \\[1mm] \textbf{Contraportada (p.8):} cierre, cita, crédito o anuncio. \\[1mm] \textbf{Sistema visual coherente:} grilla + tipografía + paleta consistentes en las 8.
\end{drawbox}

\begin{softbox}{milcMostaza}{milcGris}{Errores comunes y su corrección}
(1) Maquetar sin planear: empiezas y a la página 4 te das cuenta que no tienes contenido para la 6. (2) Variar el sistema visual por aburrimiento: cambias tipografía en la página 4 y rompes coherencia. (3) Hacer las 8 páginas iguales: sin variación de ritmo. (4) Olvidar la contraportada: la dejas en blanco o con lorem ipsum. (5) Exportar como JPG (deforma) en vez de PDF.
\end{softbox}

\begin{infoband}{milcTurquesa}


\textbf{Lo que va al cuaderno (Actividad 2 · ANALIZA).} Título: «Actividad 2 --- Plan de mi revista de 8 páginas». \textbf{Formato:} tabla de 3 columnas (Página~|~Contenido~|~Tipo de visual --- foto, texto, infografía, cita), 8 filas. \textbf{Sabes que terminaste cuando} tienes plan completo antes de tocar la herramienta digital.
\end{infoband}

\puente{Tienes el método. Toca producir: tu revista completa. Tema propio elegido por ti (cultura local, ciencia, deporte, música, etc.), 8 páginas, sistema visual coherente, exportada como PDF, publicada con URL.}

\milcestacion{milcMagenta}{Estación 3 de 3 · Hacer}

\begin{softbox}{milcMagenta}{milcGris}{Cómo vas a construir tu producto}
\textbf{Actividad 3 · APLICA --- Maqueta, exporta y publica tu revista} (40 min · individual). Esta sesión es más extensa de lo normal porque produces el producto integrador del periodo. Vas a tener tu revista de 8 páginas publicada antes de salir de aula.
\end{softbox}

\milcpaso{1}{milcMagenta}{Tu maquetación se descompone en 5 fases: (1) plan en cuaderno (8 páginas con contenido y tipo de visual), (2) maquetar en la herramienta (4 horas reales pero hoy 40 min para versión 1.0), (3) revisar coherencia visual entre las 8, (4) exportar como PDF, (5) subir a plataforma para URL pública.}
\milcpaso{2}{milcMagenta}{Sigue el patrón \textbf{``versión 1 publicada vence a versión 5 perfecta sin publicar''}: tu primera revista no será perfecta. La pregunta es si la publicas o si te quedas perfeccionándola para siempre. Lección de Periodo 1: empezar imperfecto siempre vence a esperar perfecto.}
\milcpaso{3}{milcMagenta}{Cada página expresa una idea principal. Si tu página 4 quiere decir 3 cosas, divídela en 3 páginas o elige una.}
\milcpaso{4}{milcMagenta}{Algoritmo: (1) escribe el plan de 8 páginas en cuaderno (10 min); (2) abre 8 lienzos en tu herramienta digital con tamaño A4 vertical (5 min); (3) aplica grilla guía y paleta a todas (5 min); (4) maqueta una a una empezando por la más fácil (20 min, ritmo apretado para versión 1.0); (5) revisa el conjunto (5 min); (6) exporta como PDF; (7) sube a Drive con permisos públicos.}

\begin{drawbox}{milcMagenta}{Checklist antes de declarar la revista publicada}
\checkbox 8 páginas maquetadas con tu sistema visual coherente. \\[1mm] \checkbox Portada con titular + imagen + número/fecha. \\[1mm] \checkbox Editorial/sumario en p.2. \\[1mm] \checkbox Contraportada en p.8 (no en blanco). \\[1mm] \checkbox Exportada como PDF (no JPG suelto). \\[1mm] \checkbox URL pública accesible (Drive con permisos, Issuu, Notion). \\[1mm] \checkbox Verifiqué la URL en pestaña incógnita.
\end{drawbox}

\begin{softbox}{milcMagenta}{milcGris}{Plantilla de trabajo}
\textbf{Tema de mi revista.} \textit{Frase corta.} \\[1mm] \textbf{Plan de 8 páginas.} \textit{Tabla con contenido y tipo de visual.} \\[1mm] \textbf{URL pública.} \textit{Link al PDF en Drive/Issuu/Notion.} \\[1mm] \textbf{Decisión más fuerte.} \textit{Qué decisión visual estás más orgulloso de.} \\[1mm] \textbf{Decisión más débil.} \textit{Qué cambiarías en una versión 2.0.}
\end{softbox}

\begin{infoband}{milcMagenta}


\textbf{Lo que va al cuaderno (Actividad 3 · APLICA).} Título: «Actividad 3 --- Mi revista de 8 páginas publicada». \textbf{Formato:} (1) plan de 8 páginas, (2) URL pública de la revista, (3) capturas de portada + 1 spread + contraportada, (4) reflexión personal de 3 renglones. \textbf{Extensión:} 2 páginas de cuaderno. \textbf{Sabes que terminaste cuando} tu URL es accesible y un compañero puede abrirla.
\end{infoband}

\puente{Lo que sigue resume qué entregas hoy: revista digital de 8 páginas publicada. El criterio principal: que se sienta ``un cuerpo'' coherente y la URL funcione para cualquiera con internet.}

\milcestacion{milcMostaza}{Producto de la sesión}
\textbf{Revista digital de 8 páginas con tema propio, exportada como PDF y publicada con URL}.
\begin{softbox}{milcMostaza}{milcGris}{Criterios de calidad}
(1) 8 páginas maquetadas con sistema visual coherente (grilla + tipografía + paleta). (2) Portada con titular + imagen + número/fecha. (3) Editorial o sumario en p.2. (4) Contraportada en p.8 (no en blanco). (5) Exportada como PDF. (6) URL pública accesible y verificada. (7) Reflexión personal sobre decisión más fuerte y más débil.
\end{softbox}

\puente{Antes de cerrar, mira tu revista publicada desde las cinco dimensiones humanas. Es la pieza editorial más sustantiva que has producido hasta hoy. Vale la pena nombrar qué significa.}

\milcestacion{milcVino}{Evaluación liberadora · cinco dimensiones}

\begin{softbox}{milcVino}{milcGris}{Sentido de la evaluación}
Evaluar no es solo revisar una nota. Es reconocer qué aprendí, qué cambió en mí, qué emociones manejé, cómo me relaciono con la ciudad y la comunidad, y qué le dejaré a quien viene detrás.
\end{softbox}

\textbf{Mi reflexión liberadora en cinco dimensiones.} Completa cada frase iniciada:

\small
\begin{tabularx}{\textwidth}{>{\bfseries\RaggedRight\arraybackslash}p{3.4cm}Y>{\RaggedRight\arraybackslash}p{4.6cm}}
\rowcolor{milcVino}\color{white}Dimensión & \color{white}\bfseries Mi respuesta (frame starter) & \color{white}\bfseries Algo que descubrí\\
1. Personal & Lo que más cambió en mí fue \linead{6.5cm} & \linead{4.2cm}\\
2. Emocional & La emoción más fuerte fue \linead{2.5cm} y la regulé \linead{3cm} & \linead{4.2cm}\\
3. Ciudadana & En democracia esto importa porque \linead{6cm} & \linead{4.2cm}\\
4. Local (Cartago) & En mi barrio sería útil para \linead{6.5cm} & \linead{4.2cm}\\
5. Intergeneracional & A mi abuela/abuelo le diría \linead{6.5cm} & \linead{4.2cm}\\
\end{tabularx}
\normalsize

\begin{infoband}{milcVino}
\textbf{Para la próxima sustentación.} Vuelve a esta tabla en seis meses. Verás que las respuestas cambian — no porque lo aprendido cambie, sino porque tú cambias. La evaluación liberadora es una conversación contigo a través del tiempo.
\end{infoband}


\puente{Y para terminar, tres voces sobre maquetar. Dussel sobre quién aparece en la revista, Marco Aurelio sobre la persistencia en oficio largo, Floridi sobre publicar como acto en la infoesfera.}

\milcestacion{milcGradeDark}{Triángulo de pensamiento · cierre filosófico}

\begin{softbox}{milcVino}{milcGris}{Dussel · lente del nosotros}
\textit{Toda revista es un acto editorial: decide quién aparece, quién no, en qué tono y con qué peso.}

{\footnotesize\itshape\color{milcNegro!70}--- formulación de esta guía, al modo de Enrique Dussel. No es una cita textual.}

Tu revista de 8 páginas es decisión política, no solo estética. ¿Aparecen voces de tu región, tu colegio, tu barrio? ¿O reproduces el formato de revistas nacionales que solo dan voz a élites urbanas? Cada decisión de portada y de contenido principal incluye o invisibiliza. Atreverte a poner voces tradicionalmente excluidas es soberanía editorial.

\textbf{¿Quién aparece en mi revista y qué dice eso de mis decisiones editoriales inconscientes? ¿Qué historias del Valle o de mi colegio merecen aparecer y no estaban?}
\end{softbox}
\linead{\linewidth}\\[1mm]
\linead{\linewidth}

\begin{softbox}{milcMostaza}{milcGris}{Estoicismo · Marco Aurelio · cuidado interior}
\textit{Persistir 4 horas en un oficio largo enseña más que 40 horas dispersas.}

{\footnotesize\itshape\color{milcNegro!70}--- formulación de esta guía, al modo de Marco Aurelio. No es una cita textual.}

Maquetar 8 páginas en una sola sesión exige persistencia que pocas tareas escolares piden. Es prueba estoica: aguantar el cansancio, mantener el criterio, no rendirse en la página 5 cuando te cansas. Esa persistencia es la habilidad profesional más subestimada --- y la que más distingue a quien sí termina proyectos de quien siempre los deja a medias.

\textbf{¿Qué proyectos he dejado a medias por cansancio o aburrimiento? ¿Qué cambiaría si entrenara la persistencia con esta revista como ejercicio?}
\end{softbox}
\linead{\linewidth}\\[1mm]
\linead{\linewidth}

\begin{softbox}{milcTurquesa}{milcGris}{Floridi · lente de la infoesfera}
\textit{Publicar es añadir un agente más a la infoesfera; cuídala como cuidarías un legado.}

{\footnotesize\itshape\color{milcNegro!70}--- formulación de esta guía, al modo de Luciano Floridi. No es una cita textual.}

Tu revista quedará en internet. Los archivos del futuro la guardarán. Alguien podría encontrarla en 10 años. Esa permanencia no debe paralizarte --- debe responsabilizarte. Cada decisión (qué publicas, qué ocultas, qué firmas) construye una versión de ti que vivirá más allá del aula.

\textbf{Si dentro de 5 años alguien encuentra mi revista, ¿la voy a defender o la voy a esconder? ¿Qué cambia en mi maquetación si la pienso para durar?}
\end{softbox}
\linead{\linewidth}\\[1mm]
\linead{\linewidth}

\begin{infoband}{milcNegro}
\textbf{Nota para el docente.} Las tres formulaciones son del autor de la guía, no citas textuales de los pensadores: recogen su lente para pensar el tema, no sus palabras. Si vas a citarlos en clase, remite a la obra original.
\end{infoband}

\milcestacion{milcVino}{Mi autoevaluación · marca una casilla por fila}

{\small
\setlength{\extrarowheight}{2.2pt}
\begin{tabularx}{\textwidth}{>{\RaggedRight\arraybackslash\bfseries}p{3.0cm}YYY}
\rowcolor{milcVino}\color{white}\bfseries Criterio & \color{white}\bfseries Lo logré & \color{white}\bfseries En proceso & \color{white}\bfseries Necesito apoyo\\[.6mm]
Comprensión del tema & \milcopcion{Explico las ideas con mis palabras.} & \milcopcion{Reconozco las ideas, falta precisión.} & \milcopcion{Necesito repasar lo básico.}\\[.8mm]
\rowcolor{milcGris}Pensamiento computacional & \milcopcion{Usé los pilares al armar mi producto.} & \milcopcion{Usé algunos pilares.} & \milcopcion{Necesito releer la estación 2.}\\[.8mm]
Aplicación y producto & \milcopcion{Mi producto cumple los criterios.} & \milcopcion{Está armado, con ajustes pendientes.} & \milcopcion{Aún está en construcción.}\\[.8mm]
\rowcolor{milcGris}Reflexión 5 dimensiones & \milcopcion{Respondí cada una con honestidad.} & \milcopcion{Respondí algunas con detalle.} & \milcopcion{Mis respuestas son muy generales.}\\[.8mm]
Triángulo de pensamiento & \milcopcion{Conecté las 3 voces con mi trabajo.} & \milcopcion{Conecté al menos una.} & \milcopcion{No las relaciono con mi proceso.}\\[.8mm]
\end{tabularx}}

\vspace{3mm}
\textbf{Hoy entendí que:}\\[1mm]
\linead{\linewidth}\\[1mm]
\linead{\linewidth}

\textbf{Mi compromiso para la próxima sesión es:}\\[1mm]
\linead{\linewidth}\\[1mm]
\linead{\linewidth}

\begin{softbox}{milcMostaza}{milcGris}{Cita de despedida · Marco Aurelio}
\textit{``El obstáculo en el camino se vuelve el camino. Lo que estorba al camino se vuelve camino.''}\\[1mm]
Cada error de hoy, bien observado, se convierte en pieza del camino que pisarás mañana.
\end{softbox}

\small
\begin{tabularx}{\textwidth}{>{\bfseries}p{3.6cm}Y}
\rowcolor{milcGradeDark!12}Nombre completo: & \linead{10cm}\\
Fecha: & \linead{4cm} \quad Curso/grupo: \linead{4cm}\\
Mi firma: & \linead{9cm}\\
\end{tabularx}
\normalsize

\begin{infoband}{milcGradeDark}
\textbf{Para la próxima semana.} Comparte la URL de tu revista con 3 personas distintas: un familiar, un compañero, un docente. Pide a cada uno UNA observación honesta. No discutas las observaciones: anótalas. Esas tres voces externas son insumo de la sustentación pública de la sesión 10. \emph{Felicitaciones por terminar el producto integrador del periodo.}
\end{infoband}

\end{document}
